
\thispagestyle{empty}
	
	\begin{figure}[ht]
	   \centering
			\includegraphics[width=5cm]{unal.png}
			\label{EscudoUNAL}
	\end{figure}
	
	\begin{center}
	\vspace{0.5cm}
	\LARGE
	Selección de Rasgos Funcionales de las comunidades \\
	de macroinvertebrados para su uso como variables \\
	indicadoras de  calidad en dos ríos de la sabana \\
	de Bogotá
	
	\vspace{2.5cm}
	\Large
	\textbf{Michael Sebastián Velandia Garavito}

	\vspace{4.5cm}
	\normalsize
	Universidad Nacional de Colombia \\
    Facultad de Ciencias, Departamento de Biología\\
    Bogotá, Colombia
	

	\vspace{.5cm}
	\normalsize
	Septiembre de 2021
	
	\vspace{1cm}
	\end{center}

\newpage	
\thispagestyle{empty}
	
	\begin{center}
	\vspace{0.8cm}
	\LARGE
	Selección de Rasgos Funcionales de las comunidades \\
	de macroinvertebrados para su uso como variables \\
	indicadoras de  calidad en dos ríos de la sabana \\
	de Bogotá
	
	\vspace{0.8cm}
	\Large
	\textbf{Michael Sebastián Velandia Garavito}
	
	\vspace{2.7cm}	
    \normalsize
	Trabajo de investigación presentado como requisito parcial para \\
	optar al título de: 
	
	\vspace{0.2cm}
	\normalsize
	\textbf{Biólogo}
    
	\vspace{1cm}
	\normalsize
	Director:\\
    Ph.D. Gabriel Antonio Pinilla Agudelo

	\vspace{1.3cm}
	\normalsize
	Línea de Investigación:\\
    Indicadores Biológicos en Ecosistemas Acuáticos\\
    Grupo de Investigación:\\
    Biodiversidad, biotecnología y conservación de ecosistemas\\ 
    (Depto. Biología)


	\vspace{1.3cm}
	\normalsize
	Universidad Nacional de Colombia\\
    Ciencias, Departamento de Biología\\
    Bogotá, Colombia\\
    2021\\

	\vspace{1.5cm}
	\end{center}

\newpage
\thispagestyle{empty}
    
\section*{\textit{AGRADECIMIENTOS}}
\thispagestyle{empty}

    \vspace{3cm}
    \textit{Agradezco al Centro de Estudios en Hidrobiología de la Universidad Manuela 
    Beltrán (CEH-UMB) y en especial al profesor Ciromar Lemus, quien me facilitó
    el acceso a las muestras de macroinvertebrados para su análisis y aportó la
    información de las variables fisicoquímicas medidas in situ. También agradezco
    a la Universidad Nacional de Colombia, a la Facultad de Ciencias y al Departamento
    de Biología  por el apoyo logístico para el uso de los equipos empleados para
    la identificación de los especímenes.}
    
    \vspace{.2cm}
    \textit{Agradezco al profesor Gabriel Pinilla, por su guia, su  gestión y 
    su acompañamiento durante el desarrollo de este proyecto.} 
    
    \vspace{.2cm}
    \textit{A mi familia y mi pareja, los pilares en mi vida.  Por su motivación,
    su apoyo y su paciecia, muchas gracias. No hubiera sido posible sin ustedes.}

\clearpage
